%%APPEND title="The least localizer satisfying reverse cube excision"
\textbf{Proposition (the single-square form and the cofibre test).}\ours
\textit{Let $W$ be a weak fundamental localizer satisfying the reverse implication for one cover by two sieves, in the form of W(7 bis) of Chapter~XII, ms.~76: if $X=U\cup V$ with $U,V$ sieves and $U\to X$ is in $W$, then $U\cap V\to V$ is in $W$. For a functor $g:A\to B$ let $M_g$ be the category made of $A$, of $B$, and of the arrows $a\to b$ given by the arrows $g(a)\to b$ of $B$, and let $K(g)$ be $M_g$ with one more object $t$, one arrow $a\to t$ for every object $a$ of $A$, and no other new arrow than $\id_t$. Then: (a) if $K(g)$ is $W$-aspherical, $g\in W$; (b) $W\supset W_{H\mathbb Z}$.}

\textit{Proof.}
The subcategories $M_g$ and $A_+=A\cup\{t\}$ of $K(g)$ are sieves, they cover $K(g)$, and they meet in $A$. The category $A_+$ has the terminal object $t$. If $K(g)$ is $W$-aspherical, $A_+\to K(g)$ is in $W$ by two-out-of-three, and the hypothesis gives $A\to M_g$ in $W$. The retraction $M_g\to B$ which is $g$ on $A$ has the inclusion of $B$ as right adjoint, so it is in $W_\infty\subset W$; its composite with $A\to M_g$ is $g$. This proves (a).

For (b) let $g$ be an integral homology equivalence. The nerve of $A_+$ is the cone on $NA$, so $BK(g)$ is the mapping cone of $Bg$. The function $\pi_0(g)$ is bijective; hence $C=K(g)$ is connected (it is the point $t$ if $A$ and $B$ are empty), and it is acyclic by the exact sequence of the cone. The space $BK(C\to e)$ is the unreduced suspension of $BC$: simply connected and acyclic, hence weakly contractible. So $K(C\to e)$ is $W_\infty$-aspherical, hence $W$-aspherical. By (a) for $C\to e$, the category $C$ is $W$-aspherical; by (a) for $g$, $g\in W$.\hfill$\square$

The proposition replaces the cube axiom (6b) of the previous argument by its single-square form. For localizers satisfying the relative axiom (3 bis), forward gluing gives the converse of (a), so that $g\in W$ if and only if $K(g)$ is $W$-aspherical, and this criterion is in turn equivalent to the cube forms (6b) and (6c); that equivalence is not reproduced here.

\textbf{The reversible localizers are not the localizers containing $W_{H\mathbb Z}$.}\ours
The localizer $W_0$ of the functors bijective on connected components contains $W_{H\mathbb Z}$ and does not satisfy the reverse implication. For the fold $g:\{0,1\}\to e$ the category $K(g)$ has the four objects $0,1,e,t$ and the arrows $0\to e$, $1\to e$, $0\to t$, $1\to t$; it is connected, so $W_0$-aspherical, while $g\notin W_0$. By (a), $W_0$ violates W(7 bis). The class of localizers with reverse excision has the least element $W_{H\mathbb Z}$ and is not closed upwards.

%%BLOCK after="The least localizer satisfying reverse cube excision"
\begin{dcomment}{development}{p27:je soupçonne que Loc (6 c)}{Reverse excision excludes the fibration axioms}
Reverse excision and the fibration axioms of Sections 4--6 pull in opposite directions. In this block a \emph{reversible localizer} is a fundamental localizer satisfying the relative axiom (3 bis) and the reverse implication for one cover by two sieves, as in the proposition just proved; part (a) of that proposition is called the cofibre test. Such a localizer is strongly saturated \citep[Proposition~4.2.4]{Cisinski2006}, and it contains $W_{H\mathbb Z}$. The statements below were not found in the sources consulted; no priority is claimed.

\textbf{Proposition 1 (no sphere is aspherical).}\ours
\textit{Let $W$ be a reversible localizer. If some category of the homotopy type of a sphere $S^n$, $n\geq0$, is $W$-aspherical, then $W=\operatorname{Fl}(\mathrm{Cat})$. In particular a reversible localizer other than $\operatorname{Fl}(\mathrm{Cat})$ does not contain the localizer $W_1$ of the functors inducing an equivalence of fundamental groupoids.}

\textit{Proof.}
Since $W\supset W_\infty$, asphericity depends only on the homotopy type. If $X$ has the type of $S^{n-1}$, $n\geq1$, then $K(X\to e)$ has the type of $S^n$; if it is aspherical, the cofibre test makes $X$ aspherical. Descending, the discrete category with two objects is $W$-aspherical. It is $K(\emptyset\to e)$, so $\emptyset\to e$ is in $W$ by the cofibre test. For a category $X$, the functor $\emptyset\to X$ is then in $W$ by (3 bis) over $X$: over each object $x$ it is $\emptyset\to X/x$, and $X/x$ is aspherical. Two-out-of-three gives $X\to e\in W$ for every $X$, and then $f\in W$ for every functor $f$. A category of the type of $S^2$ is $W_1$-aspherical; this gives the last assertion.\hfill$\square$

\textbf{Lemma (aspherical objects do not see free groups).}\ours
\textit{Let $W\neq\operatorname{Fl}(\mathrm{Cat})$ be a strongly saturated fundamental localizer containing $W_\infty$ for which no category of the type of $S^0$ or $S^1$ is aspherical, and let $F$ be a free group. For every $W$-aspherical category $A$ the space $BA$ is nonempty and connected, every homomorphism $\pi_1(BA)\to F$ is trivial, and the inclusion of the constant maps $BF\to\operatorname{Map}(BA,BF)$ is a weak equivalence.}

\textit{Proof.}
If $A$ were empty, $\emptyset\to e$ would be in $W$ and $W$ would be everything, as in the preceding proof. If $A$ were not connected, a surjection of $A$ onto the discrete category $\{0,1\}$ would have a section, the arrow $\{0,1\}\to e$ would be a retract of $A\to e$, and $S^0$ would be aspherical. Let $\varphi:\pi_1(BA)\to F$ be a nontrivial homomorphism. Its image is a nontrivial free group and maps onto $\mathbf Z$; a loop of $BA$ mapping to a generator and the map $BA\to S^1$ classifying the composite $\pi_1(BA)\to\mathbf Z$ make $S^1$ a retract of $BA$ in $\mathrm{Hot}$. In the localization by $W$ the image of $S^1\to e$ is then a retract of an isomorphism, hence an isomorphism, and strong saturation makes $S^1$ aspherical. So $\varphi$ is trivial. The components of $\operatorname{Map}(BA,BF)$ correspond to the conjugacy classes of homomorphisms $\pi_1(BA)\to F$, and the component of the trivial one is a $K(F,1)$ onto which the constants map by an equivalence.\hfill$\square$

\textbf{Proposition 2 (reverse excision against enough fibrations).}\ours
\textit{Let $W$ be a fundamental localizer satisfying (3 bis) which has enough intrinsic fibrations, Loc (7 F), and whose aspherical objects satisfy the conclusion of the lemma for the free group $F$ on two generators $x,y$. Then the functor $B\alpha:BF\to BF$ defined by the endomorphism $\alpha(x)=x$, $\alpha(y)=yxyx^{-1}y^{-1}$ is not in $W$. Consequently a reversible localizer with enough intrinsic fibrations is $\operatorname{Fl}(\mathrm{Cat})$.}

\textit{Proof.}
Let $W'$ be the class of functors $f$ for which $\operatorname{Map}(Bf,BF)$ is a weak equivalence. It is weakly saturated and contains $W_\infty$; it satisfies the relative Theorem~A, because $BA$ is the homotopy colimit of the spaces $B(A/c)$ over the base and $\operatorname{Map}(-,BF)$ turns homotopy colimits into homotopy limits. By hypothesis $A\to e$ is in $W'$ for every $W$-aspherical $A$.

We show $W\subset W'$. The class $\mathrm{Fib}_W$ is stable under composition and base change, every $X\to e$ is in it because $W$ is stable under products, and by (7 F) every functor is a $W$-equivalence followed by a member of $\mathrm{Fib}_W$. Brown's factorization lemma (commentary at ms.~21) then writes $f\in W$ as $f=ps$ with $p\in\mathrm{Fib}_W$ and $s$ a section of some $q\in\mathrm{Fib}_W\cap W$; so $p\in W$ too. A member $p:X\to Y$ of $\mathrm{Fib}_W\cap W$ stays in $W$ after every base change $Y'\to Y$: factor $Y'\to Y$ as $j:Y'\to Z$ in $W$ followed by $r:Z\to Y$ in $\mathrm{Fib}_W$; the pullback $p_Z$ of $p$ along $r$ is in $W$ because $r$ is an intrinsic fibration and $p\in W$; the map $X\times_YY'\to X\times_YZ$ is in $W$ because $p$ is an intrinsic fibration and $j\in W$; two-out-of-three gives the pullback of $p$ to $Y'$ in $W$. In particular $X/y\to Y/y$ is in $W$ for every object $y$, so $X/y$ is $W$-aspherical, so $X/y\to e$ is in $W'$, and Theorem~A for $W'$ gives $p\in W'$. The same holds for $q$; then $s\in W'$ and $f\in W'$.

The functor $B\alpha$ is not in $W'$. The components of $\operatorname{Map}(BF,BF)$ are the conjugacy classes of endomorphisms of $F$, and $B\alpha$ acts on them by $[\varphi]\mapsto[\varphi\alpha]$. If this were bijective there would be $\beta$ with $\beta\alpha$ inner; then $[\alpha\beta\alpha]=[\alpha]$, so $[\alpha\beta]=[\id]$ by injectivity, and $\alpha$ would be surjective. It is not: the homomorphisms $F\to S_3$ with $x\mapsto(2\,3)$ and $y\mapsto(2\,3)$, respectively $y\mapsto(1\,2)$, agree on $x$ and on $\alpha(y)$ (both send it to $(2\,3)$, since $(1\,2)(2\,3)$ has order three) and differ on $y$.

For the last assertion let $W\ne\operatorname{Fl}(\mathrm{Cat})$ be reversible. By Proposition~1 no sphere is $W$-aspherical, so the lemma applies. The endomorphism $\alpha$ is the identity on $F^{\mathrm{ab}}$, so $B\alpha$ is an integral homology equivalence and lies in $W_{H\mathbb Z}\subset W$. By the first part $W$ cannot satisfy (7 F).\hfill$\square$

\textbf{Consequences.}
(1) A reversible localizer other than $\operatorname{Fl}(\mathrm{Cat})$ contains $W_{H\mathbb Z}$, contains no truncation localizer $\mathbf W_n$, $n\geq0$ (a sphere $S^{n+1}$ is $\mathbf W_n$-aspherical), and never has enough intrinsic fibrations. It is moreover contained in $W_0$, that is, it satisfies Loc (5); this uses the description of the fundamental localizers not contained in $W_0$ and is not reproduced here.
(2) The fibration axiom Loc (7), in the form $D_{W_\omega}$, holds only for $\mathbf W_\infty$ and $\operatorname{Fl}(\mathrm{Cat})$ (commentary at ms.~65), and $\mathbf W_\infty$ is not reversible. So $\operatorname{Fl}(\mathrm{Cat})$ is the only localizer satisfying Loc (7) and reverse excision, and with Loc (5) there is none. The two groups of axioms of this chapter thus describe two different theories: ordinary homotopy types on the side of fibrations, and on the side of reversible cofibrations a family of localizers that begins at integral homology.
\end{dcomment}

%%APPEND title="The weak minimality theorem settles the finer-localizer hopes"
There is a third reading, in which the bound $\widetilde{\mathbf W}_\infty\subset\mathbf W_\infty$ is dropped and one intersects all fundamental localizers satisfying the literal list, reverse condition (6b) included. By the dichotomy just recalled, Loc~(7) leaves only $\mathbf W_\infty$ and $\operatorname{Fl}(\mathrm{Cat})$, and $\mathbf W_\infty$ fails (6b). The intersection is then $\operatorname{Fl}(\mathrm{Cat})$. In each of the three readings the conjecture $\widetilde{\mathbf W}_\infty\subsetneq\mathbf W_\infty$ is false: the minimum is $\mathbf W_\infty$, or does not exist, or is the class of all functors.

%%APPEND title="Generalized Thomason models and the role of properness"
Outside the truncations the answer can be negative. Enough intrinsic fibrations force a localizer to be generated by its aspherical objects: under Loc~(7F) every $W$-equivalence lies in every fundamental localizer containing the functors $A\to e$ with $A$ $W$-aspherical (the argument is the middle step of the proof in the next commentary). For an accessible localizer this is in turn equivalent to properness, hence to right properness of its Thomason model structure; of that equivalence only the implication from aspherical generators to properness \citep[Corollaire~4.3.28]{Cisinski2006} is used in these pages, and the other steps were not checked for this commentary. The homology localizers fail the necessary condition, as the next commentary shows; so their Thomason model structures exist, and their fibrations are not intrinsic.

%%BLOCK after="Generalized Thomason models and the role of properness"
\begin{dcomment}{development}{p124}{No homology localizer has enough intrinsic fibrations}
The question of this page for homological equivalences has a negative answer, for every generalized homology theory. For a spectrum $E$ let $W_E$ be the class of functors $f$ for which $E_*(Bf)$ is an isomorphism; for the Eilenberg--Mac~Lane spectrum of a ring $R$ this is the localizer $W_R$ of the $R$-homology equivalences. The statement was not found in the sources consulted; no priority is claimed.

\textbf{Theorem.}\ours
\textit{For every spectrum $E\neq0$ the fundamental localizer $W_E$ does not satisfy Loc (7 F): some functor has no factorization as a $W_E$-equivalence followed by an intrinsic $W_E$-fibration. In particular this holds for $W_R$, $R$ a nonzero ring.}

\textit{Proof.}
We apply the lemma and Proposition~2 of the commentary at ms.~27 (``Reverse excision excludes the fibration axioms''). The class $W_E$ is a strongly saturated fundamental localizer satisfying (3 bis) and containing $W_\infty$, and it is not $\operatorname{Fl}(\mathrm{Cat})$: $E_*(\emptyset)=0\neq E_*(e)$. The sphere $S^n$, $n\geq0$, is not $W_E$-aspherical, since $E_*(S^n)=E_*\oplus E_{*-n}$ and the projection to $E_*$ has the nonzero kernel $E_{*-n}$. So the lemma applies to $W_E$.

Let $F$ be the free group on $x,y$ and $\alpha(x)=x$, $\alpha(y)=yxyx^{-1}y^{-1}$. The space $BF$ is a wedge of two circles and $B\alpha$ induces the identity on $H_1$. Its suspension is a homology equivalence of the simply connected space $S^2\vee S^2$, hence a homotopy equivalence; by the suspension isomorphism $E_*(B\alpha)$ is an isomorphism. So $B\alpha\in W_E$. If $W_E$ satisfied (7 F), Proposition~2 would give $B\alpha\notin W_E$.\hfill$\square$

Three remarks. First, the localizers $W_E$ satisfy the cube rules (6a), (6b), (6c), by the long exact sequences of cofibre sequences; so the theorem is also an instance of the last assertion of Proposition~2. Second, $W_{H\mathbb Z}$ satisfies Loc (6a) and Loc (8) (homology commutes with filtered colimits), and it is accessible, Loc (9); the accessibility is standard for homology localizations and is not proved here. So these axioms do not imply Loc (7 F), which answers the question of ms.~105--106. Third, the truncation localizers $\mathbf W_n$ behave in the opposite way: they have enough intrinsic fibrations (preceding commentary) and fail reverse excision (commentary at ms.~17$'$).
\end{dcomment}
